\section{Related Work}

\subsection{Traditional Light Field Depth Estimation Methods}

Many depth estimation methods for light field cameras have been proposed, primarily leveraging the geometric properties of epipolar plane image (EPI)s (EPIs) [15]. Wanner and Goldluecke [16] proposed a globally consistent depth labeling framework by analyzing the dominant orientation of EPI structures using structure tensors. Specifically, the slope of the linear structures in an EPI is explicitly inversely proportional to the scene depth. Building upon this geometric foundation, Kim et al. [17] introduced an adaptive thresholding technique to refine the slope estimation, successfully improving depth accuracy in textureless regions. These EPI-based approaches demonstrate notable performance under ideal imaging conditions by exploiting the strict photo-consistency across angular views.

While EPI-based methods predominantly focus on angular correspondence, Tao et al. [18] extended the cue integration by combining defocus and correspondence cues using light field cameras. In their approach, the depth is derived by minimizing a joint energy function that aggregates both the defocus cue (measuring the sharpness of refocused images) and the correspondence cue (evaluating pixel similarity across sub-aperture images). This dual-cue fusion effectively addresses the ambiguity in highly textured regions where pure correspondence might fail. Furthermore, Lin et al. [19] refined this paradigm by introducing a robust focus measure, significantly enhancing the depth estimation quality near depth discontinuities.

To further tackle the challenges posed by occlusions and complex scene geometries, traditional works have adapted multi-view stereo (MVS) techniques to the light field domain. Jeon et al. [20] proposed an accurate depth map estimation method by employing a multi-view stereo matching framework with a phase-based sub-pixel shift calculation, which significantly improves the depth resolution. In contrast, Williem et al. [21] introduced an occlusion-noise aware data cost aggregation strategy. By explicitly modeling the occlusion boundaries and penalizing inconsistent angular views, their method successfully preserves sharp depth edges and mitigates the bleeding effects commonly observed in cost-volume based approaches.

Despite their extensive success in controlled environments, dealing with non-Lambertian surfaces remains a highly challenging problem for any kind of existing traditional light field approach [22]. First, these methods heavily rely on the Lambertian assumption, which mandates strict appearance consistency across viewpoints. Unfortunately, non-Lambertian materials, such as specular and glossy surfaces, explicitly violate this physical principle. The viewpoint-dependent reflections destroy the linear geometric constraints in EPIs, forming X-shaped patterns that lead to severe depth degradation [23]. Second, to separate specularities, some prior works attempted traditional frequency domain analyses, such as 2D Discrete Fourier Transform (2D-DFT) [8]. However, due to the inherently low angular resolution of discrete sampling (e.g., $9 \times 9$ views in commercial plenoptic cameras), the frequency resolution is insufficient to effectively decouple the ambient illumination and material reflection signals. This limitation prevents accurate physical layer decoupling. To address these limitations, we propose a novel framework which will be detailed in Section 3.

\subsection{Deep Learning-based General Light Field Depth Estimation Methods}

Driven by the success of convolutional neural networks (CNNs) in computer vision, deep learning-based methods have dramatically advanced light field depth estimation by implicitly learning the complex mapping from angular signals to depth maps [24]. Among these, EPI-based CNNs explicitly treat epipolar plane image (EPI)s as the primary input to extract angular features. For instance, Shin et al. [25] proposed EPINET, which employs 1D and 2D convolutions to extract features from horizontal and vertical EPIs, subsequently fusing them to regress depth. By leveraging the powerful representation capacity of CNNs, this approach successfully mitigates the noise sensitivity inherent in traditional structure tensor methods. However, because EPINET purely relies on data-driven feature extraction without incorporating explicit physical constraints, it implicitly assumes that the linear structures in EPIs are strictly preserved. Consequently, when dealing with non-Lambertian surfaces where EPI lines are severely distorted into $X$-shaped or curved patterns, the learned filters fail to extract valid geometric cues, leading to considerable performance degradation.

While EPI-based methods predominantly focus on angular line structures, Sub-Aperture Image (SAI) and macropixel-based approaches extend the representation space by jointly modeling spatial textures and angular parallax. A representative SAI-based work is DistgDepth proposed by Wang et al. [26], which explicitly disentangles spatial and angular information using specialized intra-view and inter-view convolutions. Specifically, the intra-view module captures rich spatial context within a single SAI, while the inter-view module aggregates parallax cues across multiple viewpoints. This decoupled architecture substantially improves depth accuracy in textureless regions by compensating for weak angular signals with strong spatial priors. Unfortunately, similar to EPI-based methods, the inter-view aggregation in DistgDepth substantially relies on the photo-consistency assumption. In mixed scenes containing glossy or specular objects, the dynamic shifting of highlights violates this assumption, causing the inter-view convolutions to aggregate conflicting features and produce severe depth hallucinations.

Alongside SAI-based methods, macropixel-based approaches rearrange the 4D light field into 2D spatial-angular grids to simultaneously capture local spatial textures and angular disparities [27]. While this representation facilitates the direct application of standard 2D CNNs and preserves local angular correlations, it inherently disrupts the global continuous angular geometry. Consequently, when applied to non-Lambertian surfaces, the irregular and viewpoint-dependent highlight shifts across macropixels introduce severe feature misalignments, substantially degrading the depth estimation performance compared to SAI-based methods.

To overcome the limited receptive fields of standard convolutions and capture global angular correlations, recent studies have introduced attention mechanisms and Transformer architectures into light field processing. Liang et al. [28] developed the Light Field Transformer (LFT), which utilizes intra-view and inter-view self-attention modules to model long-range spatial and angular dependencies. By computing attention weights across all SAIs, LFT successfully addresses the ambiguity in occluded regions by dynamically focusing on reliable, unoccluded viewpoints. Moreover, the global attention mechanism allows the network to integrate context over a much larger spatial extent. Despite these advantages, the self-attention mechanism inherently computes similarity based on feature intensity and texture. For non-Lambertian materials, the intense and viewpoint-dependent specular reflections often dominate the attention maps, misleading the network to assign high weights to physically inconsistent angular correspondences. Therefore, these Transformer-based models still struggle to establish correct angular correspondences on specular surfaces, resulting in erroneous depth predictions.

In summary, traditional methods are constrained by strict physical assumptions and limited angular resolution, while deep learning-based approaches predominantly rely on implicit data-driven mappings that fail to explicitly model the physical formation of non-Lambertian reflections. Both paradigms struggle to effectively decouple the complex angular signals caused by specularities and glossiness. To bridge this gap, our proposed method explicitly integrates physical priors into a deep learning framework, enabling robust signal decomposition and accurate depth estimation for non-Lambertian light fields, which will be elaborated in the following section.